\chapter{طرح پیشنهادی برای پیاده‌سازی در تهران}
%=====================================================================

\section{اصول طراحی}

با توجه به تجربهٔ بین‌المللی مرورشده در فصل قبل، طرح پیشنهادی برای تهران بر اصول زیر استوار است:

\begin{itemize}
\item \textbf{شروع کلان‌شهری، نه ملی:} به‌جای یک کریدور بین‌شهری پرهزینه، فاز نخست بر یک حلقهٔ \lr{QKD} کلان‌شهری با شعاع کمتر از $50$ کیلومتر متمرکز است -- مشابه رویکرد پایلوت‌های موفق \lr{HSBC} در لندن و \lr{JPMorgan} در نیویورک/نیوجرسی؛
\item \textbf{اولویت با گره‌های با بیشترین ارزش دارایی اطلاعاتی:} بانک مرکزی، مراکز دادهٔ بانک‌های بزرگ دولتی و خصوصی، و مراکز حساس دولتی؛
\item \textbf{توپولوژی حلقه‌ای \lr{(Ring)}} برای تاب‌آوری در برابر قطعی یک بخش از فیبر؛
\item \textbf{معماری ترکیبی از ابتدا:} پیاده‌سازی هم‌زمان لایهٔ \lr{PQC (ML-KEM/ML-DSA)} در نرم‌افزار، به‌موازات ساخت زیرساخت فیزیکی \lr{QKD}، تا از همان فاز نخست پوشش امنیتی فراهم باشد؛
\item \textbf{استفاده از فیبر تاریک موجود شرکت مخابرات ایران} در حد امکان، برای کاهش هزینهٔ حفاری و کابل‌کشی جدید.
\end{itemize}

\section{توپولوژی پیشنهادی}

شکل \ref{fig:tehran} توپولوژی حلقه‌ای پیشنهادی را نشان می‌دهد. شش گرهٔ اصلی (نمادین و برای اهداف نمایش فنی) در نظر گرفته شده‌اند:

\begin{enumerate}
\item \textbf{بانک مرکزی جمهوری اسلامی ایران} -- به‌عنوان گرهٔ ریشهٔ اعتماد \lr{(Root of Trust)} و مرکز مدیریت کلید سراسری؛
\item \textbf{مرکز دادهٔ اصلی شبکهٔ بانکی (شتاب/شاپرک)} -- برای ایمن‌سازی زیرساخت تسویهٔ بین‌بانکی؛
\item \textbf{مرکز دادهٔ پشتیبان بانکی} -- جهت تداوم کسب‌وکار و مسیر جایگزین؛
\item \textbf{شرکت ارتباطات زیرساخت / مخابرات ایران} -- به‌عنوان مالک و بهره‌بردار فیبر تاریک؛
\item \textbf{مرکز دادهٔ دولتی حساس (مطابق سیاست ملی فضای مجازی)}؛
\item \textbf{مرکز پژوهشی/دانشگاهی} (برای انتقال دانش فنی و پایش مستمر امنیت پروتکل).
\end{enumerate}

\begin{figure}[htbp]
\centering
\begin{tikzpicture}[
  nd/.style={rectangle,draw,thick,rounded corners,minimum width=3cm,minimum height=1.1cm,align=center,fill=green!10,font=\scriptsize},
  link/.style={thick,blue!50!black}
]
\def\R{3.6}
\def\N{6}
\node[nd,fill=goldacc!25] (r0) at (90:\R) {بانک مرکزی\\(ریشهٔ اعتماد)};
\node[nd] (r1) at (30:\R) {مرکز دادهٔ\\شتاب/شاپرک};
\node[nd] (r2) at (-30:\R) {مرکز دادهٔ\\پشتیبان بانکی};
\node[nd] (r3) at (-90:\R) {مخابرات ایران\\(هاب فیبر)};
\node[nd] (r4) at (-150:\R) {مرکز دادهٔ\\دولتی حساس};
\node[nd] (r5) at (150:\R) {مرکز\\پژوهشی/دانشگاهی};
\foreach \a/\b in {r0/r1,r1/r2,r2/r3,r3/r4,r4/r5,r5/r0}{
  \draw[link] (\a) -- (\b);
}
\node[align=center,font=\tiny] at (0,0) {توپولوژی حلقه‌ای\\(تاب‌آور در برابر\\قطعی تک‌نقطه‌ای)};
\end{tikzpicture}
\caption{توپولوژی پیشنهادی حلقهٔ \lr{QKD} کلان‌شهری تهران (فاز نخست، نمادین)}
\label{fig:tehran}
\end{figure}

\textbf{تذکر مهم:} فاصلهٔ واقعی و مکان دقیق هر یک از این نهادها باید در فاز مطالعهٔ امکان‌سنجی (فصل \lr{12}، فاز \lr{1}) با همکاری شرکت مخابرات ایران و اپراتورهای دارای مجوز فیبر تاریک، به‌طور دقیق نقشه‌برداری و تلفات واقعی فیبر اندازه‌گیری شود. با توجه به وسعت جغرافیایی تهران بزرگ، محتمل است حداقل یک گرهٔ مورد اعتماد میانی (مطابق بخش \lr{3.2}) برای اتصال گره‌های دورتر (مثلاً مراکز داده در جنوب یا غرب تهران) لازم باشد.

\section{فناوری‌های پیشنهادی برای هر لایه}

جدول \ref{tab:techstack} پشتهٔ فناوری پیشنهادی برای هر لایهٔ معماری را ارائه می‌دهد.

\begin{table}[htbp]
\centering
\caption{پشتهٔ فناوری پیشنهادی برای طرح تهران}
\label{tab:techstack}
\small
\begin{tabular}{p{3cm}p{5.8cm}p{5.5cm}}
\toprule
\textbf{لایه} & \textbf{فناوری/پروتکل پیشنهادی} & \textbf{دلیل انتخاب} \\
\midrule
لایهٔ فیزیکی \lr{QKD} & \lr{DV-QKD} مبتنی بر \lr{decoy-state BB84} با گزینهٔ ارتقا به \lr{MDI-QKD} برای گره‌های حساس & بالاترین بلوغ تجاری و بیشترین تعداد تولیدکنندهٔ معتبر (جدول \ref{tab:protocols}) \\
کانال انتقال & فیبر تک‌مود اختصاصی (تاریک) در باند \lr{O} یا \lr{C}؛ در صورت کمبود فیبر، سامانهٔ \lr{Multiplexed QKD} برای هم‌زیستی با ترافیک کلاسیک & کاهش هزینهٔ زیرساخت با اجتناب از حفاری مضاعف \Cite{toshiba2026qkd} \\
مدیریت کلید & \lr{KMS} سازگار با \lr{ETSI GS QKD 014} & سازگاری چندفروشنده‌ای و امکان ارتقای آتی \\
لایهٔ نرم‌افزاری \lr{PQC} & \lr{ML-KEM-768} برای تبادل کلید، \lr{ML-DSA-65} برای امضای تراکنش‌ها & تعادل مناسب امنیت/کارایی برای بار تراکنشی بانکی سنگین \\
رمزگذار خط \lr{(Line Encryptor)} & رمزگذار سازگار با \lr{QKD} در لایهٔ $1$/$2$ با پشتیبانی \lr{AES-256}\LTRfootnote{Advanced Encryption Standard, 256-bit key} & استاندارد صنعتی پذیرفته‌شده در پروژه‌های \lr{HSBC} و \lr{JPMorgan} \\
\bottomrule
\end{tabular}
\end{table}

%=====================================================================
